\section{Static Obfuscation}
\label{sec:obfuscation}

% --- scope-framing (short role-framing) ---
Static obfuscation transforms the data and/or weights so the server sees only
values it cannot interpret, with no---or minimal---trusted hardware: token
substitution (replacing sensitive tokens, e.g.\ SANTEXT/RANTEXT), learned embedding
transforms, covariant (data+weight) masking, null-space projection.
It is the cheapest family in our taxonomy (near-plaintext cost, no trusted hardware---Tier~1),
which is precisely why its privacy claims deserve the most scrutiny: the obfuscation
is \emph{fixed}, so a white-box adversary can analyze it in advance. We make that scrutiny
concrete in \cref{app:validation}, and the published inversions in
\cref{sec:synthesis} show several static schemes recover far more plaintext than
their authors report. On the deployment-readiness criteria,
obfuscation leads on performance and fidelity, but \emph{threat-model fit} is its weak
axis: because the transform is fixed, a white-box adversary can analyze it in advance, so its
real protection is often weaker than reported (\cref{sec:synthesis}). A second deployment
advantage is infrastructural: because the model runs unmodified, static obfuscation slots
into production inference stacks (vLLM, llama.cpp) directly---unlike crypto or TEE schemes,
which require a bespoke serving path.

\subsection{Approaches}
The defining design choice is \emph{what} the transform moves. \textbf{Data-only}
obfuscation perturbs the client's tokens or embeddings and leaves the model untouched:
token substitution (SANTEXT/CUSTEXT lineage), learned embedding transforms
(SGT~\cite{roberts_sgt}, an input-conditioned affine/Gaussian map; OSNIP~\cite{cao_osnip}, a
semantic null-space projection), and inversion-hardened projectors (Eguard~\cite{eguard},
trained against Vec2Text). Data-only transforms are attractive because the server runs a
stock model, but the model weights still encode the inverse map, so a white-box server can
attempt to undo the obfuscation---exactly the foothold the inversion attacks of
\cref{sec:obfusc-attacks} exploit.

\textbf{Covariant} obfuscation answers this by jointly transforming the data \emph{and} the
weights, so that the server computes blindly on values whose inverse it does not hold.
AloePri~\cite{lin_aloepri} applies a secret token-level permutation to the embedding and
head and absorbs it into the remaining weights via invertible matrices plus Gaussian noise;
ObfusLM~\cite{obfuslm} permutes and re-synthesizes the vocabulary, embedding, and head under
a $(k,\varepsilon)$-anonymity guarantee, the only surveyed scheme that protects the
\emph{generated output} as well as the input; and Equivariant Encryption
(EE)~\cite{buban_ee} replaces each layer's operators with transformed equivalents that
preserve the linear maps and a prescribed set of non-linearities. A third locus is the
\emph{internal state}: KV-Cloak~\cite{luo_shadowcache} obfuscates the KV-cache with a
reversible matrix and a one-time permutation refreshed per attention block, with the secret
matrices held in a TEE for key custody only (no compute is offloaded to the enclave).
Output-side and domain-specific variants round out the family: IOI~\cite{ioi} obfuscates the
\emph{decision} of a black-box classifier by mixing the instance with dummies, and
CodeCipher~\cite{codecipher} learns a token-to-token confusion map specialized to source-code
LLMs.

\subsection{Representative Schemes}
The static-obfuscation \emph{inference} schemes appear in the shared inference table
(\cref{tab:inf-perf}): the data-only transforms SGT, OSNIP, and Eguard; and the covariant /
internal-state designs AloePri, ObfusLM, KV-Cloak, and EE. All run at near-plaintext cost;
their (contested) privacy claims are examined in \cref{sec:obfusc-attacks} and
\cref{app:validation}. IOI is catalogued in prose rather than the table: it is the
decision-privacy pioneer but is confined to classification and pays an $O(mn)$ inference-call
multiplier. EE is included with a caveat---it reports near-zero overhead and exact fidelity
but offers \emph{no formal security argument} beyond the size of its key space, and is a
commercial white paper rather than a peer-reviewed result; we therefore treat its security
claim as unverified. ARoG~\cite{ning_arog}, a KGQA entity$\to$machine-ID anonymization, is a
graph-RAG defense catalogued with the retrieval schemes (\cref{tab:crypto-graph}).

A cautionary lineage is cited but excluded from the comparison: schemes whose secret is a
\emph{static} permutation or key collapse once an attacker has the public weights.
STIP~\citep{stip} (a static weight permutation), Centaur~\citep{centaur} (a static
permutation replacing the MPC non-linearities), and KV-Shield~\citep{kvshield} (a fixed
KV-cache shuffle) are each broken---by Hidden-No-More or the KV-cache Collision attack---for
the same structural reason (\cref{sec:obfusc-attacks}).

\subsection{Performance (Reported)}
\label{sec:obfusc-perf}
Online inference cost is unremarkable: every surveyed scheme runs at ${\approx}1\times$
plaintext (\cref{tab:inf-perf})---KV-Cloak adds $<$0.45\% to prefill and is lossless, EE and
AloePri are near-zero, and even the learned transforms add only milliseconds at the client.
The deployment-relevant cost is therefore \emph{offline}, and it splits the family in two.
The covariant matrix schemes (AloePri, KV-Cloak) and EE need only a one-time weight transform,
computed in minutes with no training. The learned data-only transforms pay a real training
bill---SGT trains a transform network (hours to days of GPU time), OSNIP and Eguard train an
encryptor/projection network ($1.6$--$3.4\times$ the base training), and ObfusLM additionally
requires fine-tuning the model on the obfuscated embeddings (server-side hours) before it can
serve. The headline ``${\approx}0$ inference overhead'' thus conceals a setup cost for half
the family. Set against this, the infrastructural advantage is genuine and is what makes the
family deployable today: because the served model is a stock checkpoint, AloePri and KV-Cloak
integrate directly with vLLM/SGLang and inherit KV-cache and paged-attention optimizations,
which crypto and TEE serving paths cannot. The low cost is precisely why the privacy claims
warrant the scrutiny of the next subsection.

\subsection{Limitations and Known Attacks}
\label{sec:obfusc-attacks}
Published attacks substantiate the field-level claim that static obfuscation
\emph{under-protects} relative to its reported guarantees. We group them by the schemes they
break (\cref{tab:obfusc-attacks}); the hybrid-split-specific attacks on offloaded weights and
activations (Game of Arrows, GELO's Gram/ICA leakage, static-basis recovery) are deferred to
\cref{sec:hybrid}.

\emph{Embedding inversion.} Data-only embedding transforms are the softest target. The
inversion attack on obfuscated embedding matrices~\citep{lin_inversion} recovers input tokens
from glide-reflection-style transforms, and Vec2Text-class reconstruction recovers text from
the perturbed vectors; AloePri's own evaluation reports several state-of-the-art embedding
obfuscations losing $>$30\% accuracy while still leaking $>$50\% of tokens. The defenses are
inversion-hardened training (Eguard) or a large anonymity set (ObfusLM's $k$).

\emph{Static-permutation collapse.} When the secret is a fixed permutation and the weights
are public, the adversary can solve for it. Hidden-No-More~\citep{hnm} reconstructs prompts
near-perfectly from the hidden states of permutation+noise schemes (STIP, Centaur), and the
KV-cache Collision attack~\citep{luo_shadowcache} defeats KV-Shield's fixed shuffle because
the shuffle preserves the statistical distribution the attack aligns against.

\emph{KV-cache reconstruction.} Externalizing the KV-cache opens three
attacks~\citep{luo_shadowcache}: \emph{Inversion} reconstructs the input directly from cached
keys/values, \emph{Collision} exploits the algebraic weight--cache correspondence (the one
that breaks KV-Shield), and \emph{Injection} parses the cache with a pristine model. KV-Cloak
answers all three with a per-block one-time permutation plus a beacon mask.

\emph{Deeper-layer false privacy.} The intuition that obfuscating only deep-layer
representations is safe is itself attackable: inversion succeeds at depth~\citep{depth_false_privacy},
so leaving early layers unprotected is not sufficient.

The throughline across these results is a single design lesson: the schemes that survive are
the ones whose hiding randomness is \emph{dynamic and never reused}---KV-Cloak's per-block
one-time permutation and AloePri's covariant transform on the static side, GELO's per-batch
re-mixing on the hybrid side (\cref{sec:hybrid}). The broken schemes (STIP, Centaur, KV-Shield)
all reuse a static permutation or key, which is what lets a white-box adversary average across
queries and solve for the secret.

% --- §5 attack table (parallels tab:hybrid-attacks / tab:crypto-attacks) ---
\begin{table*}[t!]
\centering
\caption{Attacks against static obfuscation. Data-only and static-permutation schemes are
broken; the surviving designs refresh their hiding randomness.}
\label{tab:obfusc-attacks}
\footnotesize
\setlength{\tabcolsep}{5pt}
\renewcommand{\arraystretch}{1.15}
\begin{tabularx}{\textwidth}{@{}>{\raggedright\arraybackslash}p{2.6cm}>{\raggedright\arraybackslash}X>{\raggedright\arraybackslash}X>{\raggedright\arraybackslash}p{3.1cm}@{}}
\toprule
\textbf{Attack} & \textbf{Target} & \textbf{Effect} & \textbf{Defense / status} \\
\midrule
Embedding inversion~\citep{lin_inversion} & data-only embedding transforms (SGT/OSNIP-class) & recovers input tokens from obfuscated embeddings; Vec2Text-class text reconstruction & inversion-hardened training (Eguard); large $k$ (ObfusLM) \\
\addlinespace[5pt]
Hidden No More~\citep{hnm} & static permutation + noise, open-weights (STIP, Centaur) & near-perfect prompt reconstruction from hidden states & dynamic per-query randomness \\
\addlinespace[5pt]
KV-cache inversion / collision / injection~\citep{luo_shadowcache} & externalized KV-cache (KV-Shield's fixed shuffle) & rebuilds the prompt from a leaked cache; Collision breaks fixed shuffles & KV-Cloak (per-block one-time perm + beacon mask) \\
\addlinespace[5pt]
Deeper-layer inversion~\citep{depth_false_privacy} & ``deep layers are safe'' assumption & inversion still succeeds at depth & protect early layers too; covariant transform \\
\bottomrule
\end{tabularx}
\end{table*}

% Validation moved to App.~\ref{app:validation}: decoupled from static-obfuscation-only and
% generalized to the heuristic-basis deployment-ready set. The §5 framing above forward-refs it.
